\section{Задача и участники}
\begin{definition}[Вариант 34: сортировочный центр посылок]
Посылки проходят входной накопитель, сканер, сортировщик, конвейер и выход своего направления. Очереди ограничены: устройство удерживает посылку до передачи дальше. Ошибки чтения вызывают повторы; после исчерпания попыток посылка поступает в ручной накопитель.
\end{definition}

C11, libc, системные вызовы, Linux. Каждая посылка, источник, сканер, сортировщик, конвейер и разгрузчик представлены процессом. Посылка выбирает повтор или переход, устройство выполняет обработку, разгрузчик освобождает накопитель. После ручного накопителя обработка не моделируется.

Число процессов с контроллером: $N+L+S+R+2D+2$, где $N$ - посылки, $L$ - источники, $S$ - сканеры, $R$ - сортировщики, $D$ - направления. По умолчанию 34 процесса.

\section{Три реализации}
\begin{center}
\begin{tabular}{>{\raggedright\arraybackslash}p{0.17\linewidth}>{\raggedright\arraybackslash}p{0.28\linewidth}>{\raggedright\arraybackslash}p{0.44\linewidth}}
\toprule
Версия & Запуск участников & Семафоры POSIX \\
\midrule
\texttt{named} & Дочерние через \texttt{fork} & Именованные, \texttt{sem\_open} \\
\texttt{unnamed} & Дочерние через \texttt{fork} & Неименованные, \texttt{sem\_init} с \texttt{pshared=1} \\
\texttt{session}, \texttt{worker} & Независимые команды & Именованные, \texttt{sem\_open} \\
\bottomrule
\end{tabular}
\end{center}

\begin{remark}[Общая память и синхронизация]
Все версии используют \texttt{shm\_open}, \texttt{ftruncate} и \texttt{mmap(MAP\_SHARED)}. В памяти находятся массивы и индексы, без указателей. Два семафора очереди считают места и элементы, бинарный семафор защищает изменения и журнал. Ожидания и задержки выполняются без блокировки. Ещё два семафора согласуют результат обработки и освобождение устройства.
\end{remark}

Независимый контроллер не создаёт участников: они подключаются по имени сеанса. Старт разрешается после регистрации всех ролей. Предел ожидания - \texttt{-{}-startup-timeout} (30 секунд). Повторные имена сеансов и роли отклоняются.

\section{Корректность и завершение}
\texttt{-{}-check} проверяет номера, единственную стадию посылки, вместимости, занятость устройств и распознанное направление. Очереди соблюдают FIFO; итоговый баланс проверяется всегда.

После доставки или отказа всех посылок возвращается 0. Обработчик SIGINT и SIGTERM сохраняет номер сигнала; рабочий код устанавливает остановку. Ожидания проверяют её до 100 мс, задержки - до 20 мс. Родитель вызывает \texttt{waitpid}; независимые участники отмечают завершение. Контроллер удаляет семафоры и shared memory. Коды: 130/143 при сигналах, 2 при неверных аргументах, 1 при ошибке.

\clearpage
\section{Интерфейс и запуск}
Аргументы задают партию, устройства, очереди, времена, ошибки, повторы и seed. Допускаются 0--512 посылок и 1--64 устройства каждого вида; справка - \texttt{-{}-help}.
\begin{lstlisting}
make
./build/center_named --check
./build/center_unnamed --parcels 17 --error-percent 100 \
    --retries 3 --check --log run.txt
python3 scripts/run_independent.py --check
make test
\end{lstlisting}

Для ручного запуска: \texttt{center\_session -{}-session demo}, затем отдельные команды \texttt{center\_worker demo ROLE INDEX}. Список ролей - \texttt{center\_worker demo plan}; полный пример приведён в README. Python-скрипт автоматизирует запуск C-программ.

\texttt{-{}-delay-ms}: 1 мс по умолчанию, 0 отключает задержки. Seed фиксирует результаты посылок, но не порядок событий. Консоль и \texttt{-{}-log} содержат PID, время, поступления, ожидания, этапы, ошибки и итоги. На Mac используется Linux-контейнер из Dockerfile.

\section{Проверки и результаты}
Среда: Docker, Linux ARM64, GCC 12.2. По журналу проверяются маршруты, FIFO, вместимости, занятость и PID. Для каждой версии проверены 25 случайных конфигураций, повторные запуски, сигналы и ошибки. Отдельно проверены частичный отказ \texttt{fork}, конфликты ролей и имён, незапущенные участники и параллельные сеансы.
\begin{center}
\begin{tabular}{>{\raggedright\arraybackslash}p{0.53\linewidth}>{\raggedright\arraybackslash}p{0.36\linewidth}}
\toprule
Сценарий для всех трёх версий & Результат \\
\midrule
40 посылок, ошибок нет & 40 доставок, 40 чтений \\
17 посылок, 100\% ошибок, 3 повтора & 17 отказов, 68 чтений \\
SIGINT/SIGTERM при ожидании & Коды 130/143, баланс сохранён \\
Штатный выход и прерывание & IPC-объекты удалены \\
\bottomrule
\end{tabular}
\end{center}
45 тестов успешны; один неприменимый тест \texttt{fork} пропущен для независимого контроллера. Журнал и инструментальные проверки - в \path{docs/DEBUGGING.md}.

\section{Исправления и ограничения}
В независимой версии итог перезаписывал начало журнала из-за отдельных файловых смещений. Исправлено открытие с \texttt{O\_APPEND}; тест сравнивает файл и консоль, включая прерывание.

Очереди работают за $O(1)$, полная проверка - за $O(N+K)$, где $K$ - число устройств и очередей. Общая блокировка ограничивает скорость, число процессов - ресурсы ОС. SIGKILL и выключение системы не гарантируют очистку.

ИИ участвовал в коде, тестах и отчёте; запросы - \path{docs/PROMPTS.md}, исправления - \path{docs/DEBUGGING.md}. Из ИДЗ2 сохранены алгоритм стадий, аргументы и оформление. Реализованы три варианта до уровня требований 7--8 баллов.
